\unit{Exposé 149 — Théorèmes de dualité pour les faisceaux algébriques cohérents}
\sid{149}
\src{169}
\seminar{(Mai 1957)}
\paperhead{THÉORÈMES DE DUALITÉ POUR LES FAISCEAUX ALGÉBRIQUES COHÉRENTS}%
  {par Alexandre GROTHENDIECK}

Les résultats qui suivent, inspirés par le \fquote{théorème de dualité algébrique}
de Serre, ont été trouvés en hiver 1955 et hiver 1956. Ils s'établissent très
simplement à l'aide de résultats assez élémentaires sur la cohomologie des espaces
projectifs [3], et l'utilisation intensive de l'algèbre homologique de
Cartan-Eilenberg, sous la forme [2].

\fsection{\normalfont1}{\normalfont\underline{Les} Ext \underline{de faisceaux de modules} ({[2]}, chap. 3 et 4).}

Soit $\symup{X}$ un espace topologique muni d'un faisceau $\symup{\Osh}$ d'anneaux avec unité
(non nécessairement commutatifs). On considère la catégorie abélienne
$\symup{\Ccat^{\Osh}}$ des faisceaux de $\symup{\Osh}$-modules, appelés aussi $\symup{\Osh}$-Modules.
On sait que tout objet de cette catégorie admet une résolution injective, ce qui
permet d'y définir les foncteurs $\symup{\operatorname{Ext}}$ ayant les propriétés
formelles bien connues. De façon précise, pour éviter des confusions, nous
désignons par $\symup{\operatorname{Hom}_{\Osh}(X;A,B)}$ ou simplement
$\symup{\operatorname{Hom}(X;A,B)}$ le \underline{groupe} abélien des $\symup{\Osh}$-homomorphismes de $\symup{A}$
dans $\symup{B}$, tandis que $\symup{\Homsh_{\Osh}(A,B)}$ désignera le \underline{faisceau} des germes
d'homomorphismes de $\symup{A}$ dans $\symup{B}$ ($\symup{A,B\in\Ccat^{\Osh}}$). On définit pour
$\symup{A\in\Ccat^{\Osh}}$ fixé des foncteurs $\symup{h_A}$ et $\symup{\hsh_A}$ à valeurs respectivement
dans la catégorie $\symup{\Ccat}$ des groupes abéliens et la catégorie
$\symup{\Ccat^{\mathrm Z}}$ des faisceaux abéliens sur $\symup{X}$, par les formules :
\begin{equation}\tag{1.1}
  \symup{h_A(B)=\operatorname{Hom}_{\Osh}(X;A,B)}
  \qquad
  \symup{\hsh_A(B)=\Homsh_{\Osh}(A,B)}
\end{equation}

Les foncteurs $\symup{h_A}$ et $\symup{\hsh_A}$ sont des foncteurs covariants exacts à gauche,
dont on considère les foncteurs dérivés droits, notés respectivement
$\symup{\operatorname{Ext}_{\Osh}^{p}(X;A,B)}$ et $\symup{\Extsh_{\Osh}^{p}(A,B)}$. On a donc
par définition
\begin{equation}\tag{1.2}
\left\{
\begin{aligned}
\symup{\operatorname{Ext}_{\Osh}^{p}(X;A,B)}
  &= \symup{(R^{p}h_A)(B)
   = H^{p}\bigl(\operatorname{Hom}_{\Osh}(X;A,C(B))\bigr)}\\
\symup{\Extsh_{\Osh}^{p}(A,B)}
  &= \symup{(R^{p}\hsh_A)(B)
   = H^{p}\bigl(\Homsh_{\Osh}(A,C(B))\bigr)}
\end{aligned}
\right.
\end{equation}

où le symbole $\symup{R^{p}}$ désigne le passage aux foncteurs dérivés droits, et où
$\symup{C(B)}$ désigne une résolution injective arbitraire de $\symup{B}$ dans $\symup{\Ccat^{\Osh}}$.
Désignons par
$\symup{\Gamma:\Ccat^{\mathrm Z}\longrightarrow\Ccat}$ le foncteur \fquote{sections}. Rappelons que ses foncteurs dérivés droits sont
notés par $\symup{B\longrightarrow H^{p}(X,B)}$ :


